% F-from-v2.tex: what Version 3 changed against Version 2, and which of its claims it withdrew.
\section{From Version 2 to Version 3}
\label{app:from-v2}

\subsection*{About this version}

Version 1 of this report (\code{polaris\_project\_report.pdf}, preserved unchanged beside this
file) described Polaris as a database-design project: twelve tables around one credential, an
entity-relationship model, a lifecycle state machine, and the argument that identity data must be
post-quantum from the first day. Version 2 (\code{polaris\_project\_report\_v2.pdf}, also preserved)
described the system as it stood at tag \code{v9.345}, thirty-six tables, a federation protocol, a
transparency layer and an institutional exchange fabric, and was revised at \code{v9.355} when the
holder-side work it had recorded as missing was built. Version 3 describes the system at version
\code{1.0.0-rc.7} (commit \code{4e74376}, 20 September 2026), re-measured whole rather than patched,
and it is written for a reader who has never opened the repository. It was first written at
\code{1.0.0-rc.1} four days earlier and re-measured in place; fourteen of the twenty-six counts in
Appendix~\ref{app:counts} did not move, and the ones that did are marked there.

\textbf{What changed, in one paragraph.} Between Version 2 and this one the repository stopped
counting its own progress and started counting witnesses. The four standalone packages a relying
party installs are on public registries. A wallet written by somebody else, run unmodified,
presented a credential to the verifier and was accepted, after refusing it once for a defect every
internal instrument had passed over. The OpenID Foundation's hosted conformance suite ran the
verifier profile across the open internet. The verification layer itself was attacked from inside:
the checks, the constraints, the triggers, the unique indexes, the stored procedures, the zero-knowledge
witnesses, the conformance contract and both reference verifiers were each mutated and required to
go red, and many did not, at first. The physical layer became an object somebody else can implement.
And two words on the front door were weakened, \emph{post-quantum} to \emph{algorithm agility under an
audited migration path} and \emph{compulsion-resistant} to \emph{duress-aware}, because the evidence
did not support the stronger ones. Appendix~\ref{app:from-v2} lists every claim of Version 2 that this
version withdraws or narrows.

\textbf{Three dates, stated plainly.} Version 2's body was measured at \code{v9.345}; its
holder-side sections at \code{v9.355}. Everything in this version, every count, every mark and every
figure, was measured at \code{1.0.0-rc.7} by the methods in Appendix~\ref{app:counts}. Where a number
is older than that, because the repository has not re-run the instrument that produced it, the paper
says so beside the number. Nothing that ran at \code{v9.355} stopped running.

The paper keeps what Version 1 got right, the token-centred geometry of the schema, used here as the
centre of every map, and the lifecycle machine, and what Version 2 got right, the concentric reading
of the system, one ring per section, each ring introduced by the failure of the ring inside it.

Version 2 is preserved unchanged beside this file. This appendix records what changed between
them, and, more usefully, which sentences of Version 2 this version can no longer write.

\subsection*{What the system gained}

Between \code{v9.355} and \code{1.0.0-rc.1}, in the order the rings meet them: what an enrolment
rests on, the trusted referee and the enrolment code; the card as an encoding, an emulator, a
personalization step and a verifier device, with duress measured in time; decisions about
authorities, operators, quotas, ceilings and relying parties kept as append-only records; eight
mutation drills over the verification layer and the corrections they forced; four formal models
checked on every push; a population migration that runs and is measured; the cache rules that let an
untrusted network carry a status; a depth-24 epoch pipeline; a measured statement of what two
colluding verifiers can do; a standby region; a capacity model that found an integer column, and a
cost model; the public transparency report; the pilot wind-down and the coexistence plan; the
machine-checked assurance mapping and the automated third of accessibility; two format bridges; four
packages on public registries; the OpenID4VP verifier; the operating contract and the scoreboard; and
the first three external nouns on it.

\subsection*{Claims of Version 2 that this version withdraws or narrows}

{\footnotesize
\setlength{\tabcolsep}{4pt}\renewcommand{\arraystretch}{1.12}
\rowcolors{2}{tabstripe}{white}
\begin{xltabular}{\textwidth}{>{\RaggedRight\arraybackslash\hspace{0pt}}p{2.0in} >{\RaggedRight\arraybackslash\hspace{0pt}}X >{\RaggedRight\arraybackslash\hspace{0pt}}p{1.6in}}
\caption{Sentences of Version 2 that Version 3 does not repeat.}\label{tab:withdrawn}\\
\toprule
\rowcolor{tabheader}\thead{Version 2 said} & \thead{Version 3 says} & \thead{Why} \\
\midrule
\endfirsthead
\toprule
\rowcolor{tabheader}\thead{Version 2 said} & \thead{Version 3 says} & \thead{Why} \\
\midrule
\endhead
\midrule\multicolumn{3}{r}{\itshape continued on the next page}\\
\endfoot
\bottomrule
\endlastfoot
A post-quantum, issuer-unlinkable, compulsion-resistant identity-token system & An issuer-unlinkable, duress-aware identity-token system signed with ML-DSA-65 under an audited algorithm-migration path & Two words weakened by the repository's own assessments; each is now refused by a check on outward surfaces \\
The default signing algorithm is post-quantum from day one & The algorithm is ML-DSA-65 and the default development and CI path signs a labelled placeholder; production fails closed without real signing & Measured: 32 placeholder bytes where a signature is 3,309 \\
The duress path is a tested property; no external party has measured its side channel & The mechanism resists the casual coercer, not the informed one beside the holder, and is net-negative against lawful access & The lab's assessment against the contract's three coercers \\
An actual implementation by an external team has not happened & One unmodified outside wallet has presented over OpenID4VP and been accepted; no implementation of the native protocol by anyone else & The scoreboard's wallet row; the boundary between the two paths \\
The specification and the conformance suite exist; the integration is unproven & The OpenID Foundation's hosted suite ran the profile: 7 passed, 4 in review; a later run against \code{1.0.0rc7} was certified (24 September 2026) & The Foundation's own signed exports; its certified-implementation listing \\
Cross-verifier correlation is a permanent, documented property, then bounded & Bounded on the zero-knowledge path within the epoch's crowd, against an adversary that matches equality; exposed on plain and SD-JWT presentations; the crowd has no floor & The linkability lab's three findings \\
The certificate the verifier signs requests with was correct & It carried no key usage, and every internal instrument had passed it & An outside wallet refused it \\
The check layer catches its own violation & Seventy-one of seventy-five checks passed on a tree where their property was commented out, and the tests for constraints, triggers, indexes and procedures would often not have noticed their subject missing & The mutation drills; every finding closed and drilled \\
Thirty-six tables, 26 triggers, 191 checks, 48 conformance cases, 18 CI jobs, 20 formats, 1,048 test functions, 755 and 84 of 89 tests & Forty-five tables, 42 triggers, 272 checks, 118 cases, 20 jobs, 27 formats, 1,755 test functions, 992 and 107 of 112 tests & Re-measured at \code{1.0.0-rc.1}; Appendix~\ref{app:counts} \\
Thirteen audit-of-record instances & Fourteen named, thirty-one tables under an append-only trigger & The schema's own triggers, counted \\
The physical token is modelled and not manufactured & Encoded, emulated, personalized and read by a reference device; not manufactured & \code{polaris\_card/} and its drills \\
Multi-region recovery is not built & A standby region is built and its non-zero recovery point is measured & The evacuation drill \\
The roadmap's targets are to be validated & A throughput model built from single-machine measurements clears every target by an order of magnitude under stated assumptions, nothing ran at national scale, and a 32-bit column would have ended the service in five days & The capacity model \\
The version is v9.355, a count of ships & The version is \code{1.0.0-rc.7}, moved only by an externally observable change & The operating contract's versioning rule \\
Future work is a list of things to build & Future work is a list of things outside parties would have to do, questions the lab owes answers to, and items held until somebody outside needs them & The operating contract's merge rule and kill line \\
\end{xltabular}
}

\subsection*{What was deliberately not done}

The dyno and the application baseline were not re-run, and their rows in Table~\ref{tab:perf}
carried their original stamps; the dyno was re-run on 23 September 2026 (Appendix~\ref{app:counts}). The Version 1 and Version 2 documents were not edited. The repository
documents that disagree with the code at \code{1.0.0-rc.7} were not corrected for the paper; they
are listed in the author's notes, as Version 2's were, so that the paper follows the code and the
disagreement is not lost.

\subsection*{Vocabulary introduced}

A fifth status mark, \sWit, for a capability a named outside party has exercised, with the rule
that it is given only where the scoreboard has a filled row and never means audited, certified or
piloted. A thirteenth rung on the ladder, for the knowledge the tree cannot check about itself. And
the operating contract's own words, external noun, scoreboard, release candidate, used the way the
repository uses them.
